\chapter{35}



Louis klickte sich durch die angezeigten Artikel. Dann, unter «Comune di
Milano». Da war es: «Polizeiliche Mitteilungen - Unfälle, Verbrechen und
Zeugenaufrufe». Seine Augen weiteten sich. Ihm wurde abwechselnd heiss
und kalt. Er schluckte und begann zu lesen.

\sectionbreak

\indentedblock{Zeugenaufruf «Roccia del Culto», Milano, 25. August

Eine 26-jährige Frau wurde am frühen Montag des 20. August auf einen
anonymen Hinweis hin, schwer verletzt in der Stadt Mailand am Fusse des
«Roccia del Culto», durch die örtliche Feuerwehr geborgen.

Ein Verbrechen kann nicht ausgeschlossen werden. Die Polizei tappt nach
wie vor im Dunkeln.

Die Polizia di Stato Milano bittet die Bevölkerung um ihre Mithilfe.

Wer hat sich in der Nacht vom 19. auf den 20. August, zwischen
Mitternacht und 04.00 Uhr rund um den «Roccia del Culto» aufgehalten,
etwas Auffälliges beobachtet oder gehört?

Sachdienliche Hinweise bitte direkt an den zuständigen Commissario
Alessandro Mantovani der Polizia di Stato Milano - hier folgte eine
Telefonnummer - oder an jede andere Polizeidienststelle.}

\sectionbreak

Phua. Nichts Neues zwar, doch diesmal schwarz auf weiss vor seinen
Augen. Und wohl auf ewig erhalten. Louis las die Nachricht aufs Neue.
Und noch einmal. Er hielt den Atem an. Bis er sich auf sein Bett fallen
liess. Seine Finger begannen augenblicklich zu zittern. Er liess das Handy über dem Bauch
fallen.
Der Zeugenaufruf löste eine Kette von Fragen aus. Sie schienen
unbeantwortet an ihm vorbeizusausen, um gleich von neuem aufzutauchen.
Er starrte an die Decke. Sein Magen begann sich zu drehen. Johs Musik
von unten fehlte ihm gerade noch. \emph{«079»}. Der
Schweizer Sommerhit des Jahres, wie Louis später wissen würde. Er kam aus
einer ihm aktuell völlig fremden Welt. Joh putzte offenbar bevorzugt mit
lauter Radiomusik.
Die melodiöse Leichtigkeit des Liedes ging Louis gerade massiv auf die Nerven. 

\qrblock{Lo \& Leduc \emph{«079»}}{media/QR-Codes/030_079_Lo_&_Leduc_Spotify.png}

\emph{Schwer verletzt, von der Feuerwehr geborgen}? Die Information
stimmte mit Mamans Aussage überein. Das hiess wenigstens, Giorgia lebte.
Aber warum, zum Teufel, musste sie von der Feuerwehr gerettet werden?
Louis stellte sich den Unfallort vor. Überlegte hin und her. Er fand
einfach keinen plausiblen Grund.
Und, wer hatte den Vorfall gemeldet? Woher kam der anonyme Hinweis? Warum
denn anonym? Warum nur hatte er sie nicht gefunden?

Wie ging es ihr? Louis graute.

Die Polizei suchte also nach Zeugen. Er wehrte sich gegen die
Vorstellung, Giorgia schwer verletzt im Spital zu wissen. Wenn
öffentlich Zeugen gesucht wurden, musste es unweigerlich darum gehen,
Schuldfragen zu klären. Worum denn sonst? Und Schuldfragen wurden nur
geklärt, wenn etwas Undurchsichtiges passiert war.
Immerhin, sein Name war nirgendwo zu lesen. Eine Suche nach ihm fand
er auch nach weiterem Googeln nicht. Aber vielleicht hatte das mit der
Strategie der Polizei zu tun? Oder die Ermittler standen noch immer bei
Null? Was war wahrscheinlicher?
Es gab keine Antworten. Er konnte nichts tun.
Mit einem Ruck stand er auf, stellte sich vor den Waschtisch und begann
sich mit beiden Händen kaltes Wasser ins Gesicht zu klatschen.

Die Polizeimeldung hatte ihn mit einem Schlag zurück in die Verzweiflung
geschleudert. 
Unten wartete Joh. Louis musste unbeschwert auftreten. Er brauchte diese
temporäre Anstellung. Dringend. Er war endgültig ausgebrannt. Gut, waren wenigstens seine
Fixkosten gedeckt. Die Daueraufträge würden vorläufig problemlos weiterlaufen.




